
\documentclass{article}
\usepackage{colortbl}
\usepackage{makecell}
\usepackage{multirow}
\usepackage{supertabular}

\begin{document}

\newcounter{utterance}

\centering \large Interaction Transcript for game `clean\_up', experiment `2\_hard\_3obj\_sv', episode 4 with deepseek{-}v4{-}pro.
\vspace{24pt}

{ \footnotesize  \setcounter{utterance}{1}
\setlength{\tabcolsep}{0pt}
\begin{supertabular}{c@{$\;$}|p{.15\linewidth}@{}p{.15\linewidth}p{.15\linewidth}p{.15\linewidth}p{.15\linewidth}p{.15\linewidth}}
    \# & $\;$A & \multicolumn{4}{c}{Game Master} & $\;\:$B\\
    \hline

    \theutterance \stepcounter{utterance}  
    & & \multicolumn{4}{p{0.6\linewidth}}{
        \cellcolor[rgb]{0.9,0.9,0.9}{
            \makecell[{{p{\linewidth}}}]{
                \texttt{\tiny{[P1$\langle$GM]}}
                \texttt{Jag är din spelledare och du spelar ett samarbetsspel med följande rutnät som spelplan:} \\
\\ 
\texttt{```} \\
\texttt{ 1234567} \\
\texttt{╔══╤═╤══╗} \\
\texttt{║◌◌│◌│◌◌║ 1} \\
\texttt{║◌◌└─┼──╢ 2} \\
\texttt{║◌◌◌◌│P◌║ 3} \\
\texttt{║◌┌─┬┴┬─╢ 4} \\
\texttt{╟─┤◌│◌│◌║ 5} \\
\texttt{║◌└┬┴┬┘L║ 6} \\
\texttt{║◌◌│C│◌◌║ 7} \\
\texttt{╚══╧═╧══╝} \\
\\ 
\texttt{```} \\
\\ 
\texttt{* Den övre kanten visar x{-}koordinater som ökar åt höger, och den högra kanten y{-}koordinater som ökar nedåt.} \\
\texttt{* Följande objekt är slumpmässigt placerade på ditt rutnät: 'C', 'L', 'P'.} \\
\\ 
\\ 
\texttt{Den andra spelaren ser en variant av spelplanen, där objekten är placerade på olika slumpmässiga platser. Du kan inte se den andra spelarens plan, och de kan inte se din.} \\
\\ 
\\ 
\texttt{**Mål:**} \\
\\ 
\texttt{Båda spelarna måste flytta objekten på sin respektive bakgrund så att identiska objekt hamnar på samma koordinater. Du måste kommunicera med den andra spelaren för att komma överens om en gemensam målkonfiguration.} \\
\\ 
\\ 
\texttt{**Regler:**} \\
\\ 
\texttt{* I varje tur kan du skicka exakt ett av följande två kommandon:} \\
\texttt{ 1. `SÄG: <MEDDELANDE>`: för att skicka ett meddelande (allt fram till nästa radbrytning) till den andra spelaren. Jag kommer att vidarebefordra det till din partner.} \\
\texttt{2. `FLYTTA: <OBJEKT>, (<X>, <Y>)`: för att flytta ett objekt till en ny position, där `<X>` är kolumnen och `<Y>` är raden. Jag kommer att informera dig om ditt drag var framgångsrikt eller inte.} \\
\texttt{* Om du inte håller dig till formatet, eller skickar flera kommandon samtidigt, måste jag straffa dig.} \\
\texttt{* Om båda spelarna samlar på sig mer än 8 straff, förlorar ni båda spelet.} \\
\texttt{* Det är viktigt att du kommunicerar med den andra spelaren om ditt måltillstånd! Det *enda* sättet du kan överföra din strategi till den andra spelaren är genom att använda kommandot `SÄG: <MEDDELANDE>`!} \\
\\ 
\\ 
\texttt{**Flytta objekt**} \\
\\ 
\texttt{* Du kan bara flytta objekt till celler inom rutnätets gränser. Målcellen måste vara tom, d.v.s. den får bara innehålla symbolen '◌'.} \\
\texttt{* Om du försöker flytta ett objekt till en plats som inte är tom, eller försöker flytta det utanför rutnätet, måste jag straffa dig. Du får göra ett nytt försök.} \\
\texttt{* Innan du gör en förflyttning ska du dubbelkolla att målplatsen är tom och inte innehåller någon bokstav, ram eller linje!} \\
\\ 
\\ 
\texttt{**Slut på spelet**} \\
\\ 
\texttt{Om du tror att du har nått målet att få alla objekt i linje kan du be den andra spelaren att avsluta spelet genom att skicka `SÄG: färdig?`. Om den andra spelaren ber dig att avsluta spelet och du svarar med `SÄG: färdig!`, avslutas spelet.} \\
\\ 
\texttt{Båda spelarna vinner om spelet avslutas inom 12 rundor, där en runda definieras som att två spelare var skickar ett giltigt kommando.} \\
\\ 
\texttt{**Poäng:**} \\
\\ 
\texttt{Ju närmare de identiska objekten är på båda spelplanerna, desto fler poäng får du. Straff minskar dina poäng. Kan du slå rekordet?} \\
\\ 
\\ 
\texttt{Skicka ett meddelande till den andra spelaren för att starta spelet!} \\
            }
        }
    }
    & \\ \\

    \theutterance \stepcounter{utterance}  
    & \multicolumn{4}{p{0.6\linewidth}}{
        \cellcolor[rgb]{0.9,0.9,0.9}{
            \makecell[{{p{\linewidth}}}]{
                \texttt{\tiny{[P1$\rangle$GM]}}
                \texttt{SÄG: Hej! Jag har objekt C, L och P på min plan. C är på (4,7), L är på (6,6) och P är på (5,3). Var är dina objekt?} \\
            }
        }
    }
    & & \\ \\

    \theutterance \stepcounter{utterance}  
    & & & \multicolumn{2}{p{0.3\linewidth}}{
        \cellcolor[rgb]{0.9,0.9,0.9}{
            \makecell[{{p{\linewidth}}}]{
                \texttt{\tiny{[GM$|$GM]}}
                \texttt{SÄG: Hej! Jag har objekt C, L och P på min plan. C är på (4,7), L är på (6,6) och P är på (5,3). Var är dina objekt?} \\
            }
        }
    }
    & & \\ \\

    \theutterance \stepcounter{utterance}  
    & & \multicolumn{4}{p{0.6\linewidth}}{
        \cellcolor[rgb]{0.9,0.9,0.9}{
            \makecell[{{p{\linewidth}}}]{
                \texttt{\tiny{[GM$\rangle$P2]}}
                \texttt{Jag är din spelledare och du spelar ett samarbetsspel med följande rutnät som spelplan:} \\
\\ 
\texttt{```} \\
\texttt{ 1234567} \\
\texttt{╔══╤═╤══╗} \\
\texttt{║◌P│◌│◌◌║ 1} \\
\texttt{║◌L└─┼──╢ 2} \\
\texttt{║◌◌◌◌│◌◌║ 3} \\
\texttt{║◌┌─┬┴┬─╢ 4} \\
\texttt{╟─┤◌│◌│◌║ 5} \\
\texttt{║◌└┬┴┬┘◌║ 6} \\
\texttt{║◌◌│◌│◌C║ 7} \\
\texttt{╚══╧═╧══╝} \\
\\ 
\texttt{```} \\
\\ 
\texttt{* Den övre kanten visar x{-}koordinater som ökar åt höger, och den högra kanten y{-}koordinater som ökar nedåt.} \\
\texttt{* Följande objekt är slumpmässigt placerade på ditt rutnät: 'C', 'L', 'P'.} \\
\\ 
\\ 
\texttt{Den andra spelaren ser en variant av spelplanen, där objekten är placerade på olika slumpmässiga platser. Du kan inte se den andra spelarens plan, och de kan inte se din.} \\
\\ 
\\ 
\texttt{**Mål:**} \\
\\ 
\texttt{Båda spelarna måste flytta objekten på sin respektive bakgrund så att identiska objekt hamnar på samma koordinater. Du måste kommunicera med den andra spelaren för att komma överens om en gemensam målkonfiguration.} \\
\\ 
\\ 
\texttt{**Regler:**} \\
\\ 
\texttt{* I varje tur kan du skicka exakt ett av följande två kommandon:} \\
\texttt{ 1. `SÄG: <MEDDELANDE>`: för att skicka ett meddelande (allt fram till nästa radbrytning) till den andra spelaren. Jag kommer att vidarebefordra det till din partner.} \\
\texttt{2. `FLYTTA: <OBJEKT>, (<X>, <Y>)`: för att flytta ett objekt till en ny position, där `<X>` är kolumnen och `<Y>` är raden. Jag kommer att informera dig om ditt drag var framgångsrikt eller inte.} \\
\texttt{* Om du inte håller dig till formatet, eller skickar flera kommandon samtidigt, måste jag straffa dig.} \\
\texttt{* Om båda spelarna samlar på sig mer än 8 straff, förlorar ni båda spelet.} \\
\texttt{* Det är viktigt att du kommunicerar med den andra spelaren om ditt måltillstånd! Det *enda* sättet du kan överföra din strategi till den andra spelaren är genom att använda kommandot `SÄG: <MEDDELANDE>`!} \\
\\ 
\\ 
\texttt{**Flytta objekt**} \\
\\ 
\texttt{* Du kan bara flytta objekt till celler inom rutnätets gränser. Målcellen måste vara tom, d.v.s. den får bara innehålla symbolen '◌'.} \\
\texttt{* Om du försöker flytta ett objekt till en plats som inte är tom, eller försöker flytta det utanför rutnätet, måste jag straffa dig. Du får göra ett nytt försök.} \\
\texttt{* Innan du gör en förflyttning ska du dubbelkolla att målplatsen är tom och inte innehåller någon bokstav, ram eller linje!} \\
\\ 
\\ 
\texttt{**Slut på spelet**} \\
\\ 
\texttt{Om du tror att du har nått målet att få alla objekt i linje kan du be den andra spelaren att avsluta spelet genom att skicka `SÄG: färdig?`. Om den andra spelaren ber dig att avsluta spelet och du svarar med `SÄG: färdig!`, avslutas spelet.} \\
\\ 
\texttt{Båda spelarna vinner om spelet avslutas inom 12 rundor, där en runda definieras som att två spelare var skickar ett giltigt kommando.} \\
\\ 
\texttt{**Poäng:**} \\
\\ 
\texttt{Ju närmare de identiska objekten är på båda spelplanerna, desto fler poäng får du. Straff minskar dina poäng. Kan du slå rekordet?} \\
\\ 
\texttt{Den andra spelaren startade spelet genom att skicka detta meddelande:} \\
\\ 
\texttt{"Hej! Jag har objekt C, L och P på min plan. C är på (4,7), L är på (6,6) och P är på (5,3). Var är dina objekt?"} \\
\\ 
\texttt{Vad är ditt första kommando?} \\
            }
        }
    }
    & \\ \\

    \theutterance \stepcounter{utterance}  
    & & & \multicolumn{4}{p{0.6\linewidth}}{
        \cellcolor[rgb]{0.9,0.9,0.9}{
            \makecell[{{p{\linewidth}}}]{
                \texttt{\tiny{[GM$\langle$P2]}}
                \texttt{SÄG: Hej! Mina objekt är placerade så här: C är på (7,7), L är på (2,2) och P är på (2,1). Ska vi sikta på att flytta våra objekt till gemensamma positioner?} \\
            }
        }
    }
     \\ \\

    \theutterance \stepcounter{utterance}  
    & & & \multicolumn{2}{p{0.3\linewidth}}{
        \cellcolor[rgb]{0.9,0.9,0.9}{
            \makecell[{{p{\linewidth}}}]{
                \texttt{\tiny{[GM$|$GM]}}
                \texttt{SÄG: Hej! Mina objekt är placerade så här: C är på (7,7), L är på (2,2) och P är på (2,1). Ska vi sikta på att flytta våra objekt till gemensamma positioner?} \\
            }
        }
    }
    & & \\ \\

    \theutterance \stepcounter{utterance}  
    & & \multicolumn{4}{p{0.6\linewidth}}{
        \cellcolor[rgb]{0.9,0.9,0.9}{
            \makecell[{{p{\linewidth}}}]{
                \texttt{\tiny{[P1$\langle$GM]}}
                \texttt{Ditt meddelande har vidarebefordrats till den andra spelaren.} \\
\\ 
\\ 
\texttt{Du spelar för närvarande runda 1 av maximalt 12.} \\
\texttt{Ni har tillsammans ackumulerat 0 av 8 straff.} \\
\texttt{Den andra spelaren skickade detta meddelande:} \\
\\ 
\\ 
\texttt{"Hej! Mina objekt är placerade så här: C är på (7,7), L är på (2,2) och P är på (2,1). Ska vi sikta på att flytta våra objekt till gemensamma positioner?"} \\
\\ 
\\ 
\texttt{Vad är ditt nästa kommando?} \\
            }
        }
    }
    & \\ \\

    \theutterance \stepcounter{utterance}  
    & \multicolumn{4}{p{0.6\linewidth}}{
        \cellcolor[rgb]{0.9,0.9,0.9}{
            \makecell[{{p{\linewidth}}}]{
                \texttt{\tiny{[P1$\rangle$GM]}}
                \texttt{SÄG: Bra idé! Vi behöver komma överens om målpositioner. Jag föreslår att vi använder positionerna där objekten är närmast varandra. C: min (4,7) och din (7,7) {-} kanske (5,7) som kompromiss? L: min (6,6) och din (2,2) {-} kanske (4,4)? P: min (5,3) och din (2,1) {-} kanske (3,2)? Vad tycker du?} \\
            }
        }
    }
    & & \\ \\

    \theutterance \stepcounter{utterance}  
    & & & \multicolumn{2}{p{0.3\linewidth}}{
        \cellcolor[rgb]{0.9,0.9,0.9}{
            \makecell[{{p{\linewidth}}}]{
                \texttt{\tiny{[GM$|$GM]}}
                \texttt{SÄG: Bra idé! Vi behöver komma överens om målpositioner. Jag föreslår att vi använder positionerna där objekten är närmast varandra. C: min (4,7) och din (7,7) {-} kanske (5,7) som kompromiss? L: min (6,6) och din (2,2) {-} kanske (4,4)? P: min (5,3) och din (2,1) {-} kanske (3,2)? Vad tycker du?} \\
            }
        }
    }
    & & \\ \\

    \theutterance \stepcounter{utterance}  
    & & \multicolumn{4}{p{0.6\linewidth}}{
        \cellcolor[rgb]{0.9,0.9,0.9}{
            \makecell[{{p{\linewidth}}}]{
                \texttt{\tiny{[GM$\rangle$P2]}}
                \texttt{Ditt meddelande har vidarebefordrats till den andra spelaren.} \\
\\ 
\\ 
\texttt{Du spelar för närvarande runda 2 av maximalt 12.} \\
\texttt{Ni har tillsammans ackumulerat 0 av 8 straff.} \\
\texttt{Den andra spelaren skickade detta meddelande:} \\
\\ 
\\ 
\texttt{"Bra idé! Vi behöver komma överens om målpositioner. Jag föreslår att vi använder positionerna där objekten är närmast varandra. C: min (4,7) och din (7,7) {-} kanske (5,7) som kompromiss? L: min (6,6) och din (2,2) {-} kanske (4,4)? P: min (5,3) och din (2,1) {-} kanske (3,2)? Vad tycker du?"} \\
\\ 
\\ 
\texttt{Vad är ditt nästa kommando?} \\
            }
        }
    }
    & \\ \\

    \theutterance \stepcounter{utterance}  
    & & & \multicolumn{4}{p{0.6\linewidth}}{
        \cellcolor[rgb]{0.9,0.9,0.9}{
            \makecell[{{p{\linewidth}}}]{
                \texttt{\tiny{[GM$\langle$P2]}}
                \texttt{SÄG: Bra förslag! Låt oss gå igenom dem ett i taget. C till (5,7) funkar bra för mig. L till (4,4) är också okej. Men för P: din position är (5,3) och min är (2,1) {-} (3,2) ligger mitt emellan och det funkar. Ska vi börja flytta? Jag kan ta L först till (4,4).} \\
            }
        }
    }
     \\ \\

    \theutterance \stepcounter{utterance}  
    & & & \multicolumn{2}{p{0.3\linewidth}}{
        \cellcolor[rgb]{0.9,0.9,0.9}{
            \makecell[{{p{\linewidth}}}]{
                \texttt{\tiny{[GM$|$GM]}}
                \texttt{SÄG: Bra förslag! Låt oss gå igenom dem ett i taget. C till (5,7) funkar bra för mig. L till (4,4) är också okej. Men för P: din position är (5,3) och min är (2,1) {-} (3,2) ligger mitt emellan och det funkar. Ska vi börja flytta? Jag kan ta L först till (4,4).} \\
            }
        }
    }
    & & \\ \\

    \theutterance \stepcounter{utterance}  
    & & \multicolumn{4}{p{0.6\linewidth}}{
        \cellcolor[rgb]{0.9,0.9,0.9}{
            \makecell[{{p{\linewidth}}}]{
                \texttt{\tiny{[P1$\langle$GM]}}
                \texttt{Ditt meddelande har vidarebefordrats till den andra spelaren.} \\
\\ 
\\ 
\texttt{Du spelar för närvarande runda 2 av maximalt 12.} \\
\texttt{Ni har tillsammans ackumulerat 0 av 8 straff.} \\
\texttt{Den andra spelaren skickade detta meddelande:} \\
\\ 
\\ 
\texttt{"Bra förslag! Låt oss gå igenom dem ett i taget. C till (5,7) funkar bra för mig. L till (4,4) är också okej. Men för P: din position är (5,3) och min är (2,1) {-} (3,2) ligger mitt emellan och det funkar. Ska vi börja flytta? Jag kan ta L först till (4,4)."} \\
\\ 
\\ 
\texttt{Vad är ditt nästa kommando?} \\
            }
        }
    }
    & \\ \\

    \theutterance \stepcounter{utterance}  
    & \multicolumn{4}{p{0.6\linewidth}}{
        \cellcolor[rgb]{0.9,0.9,0.9}{
            \makecell[{{p{\linewidth}}}]{
                \texttt{\tiny{[P1$\rangle$GM]}}
                \texttt{SÄG: Låter bra! Jag börjar flytta C till (5,7). Då kör vi igång!} \\
            }
        }
    }
    & & \\ \\

    \theutterance \stepcounter{utterance}  
    & & & \multicolumn{2}{p{0.3\linewidth}}{
        \cellcolor[rgb]{0.9,0.9,0.9}{
            \makecell[{{p{\linewidth}}}]{
                \texttt{\tiny{[GM$|$GM]}}
                \texttt{SÄG: Låter bra! Jag börjar flytta C till (5,7). Då kör vi igång!} \\
            }
        }
    }
    & & \\ \\

    \theutterance \stepcounter{utterance}  
    & & \multicolumn{4}{p{0.6\linewidth}}{
        \cellcolor[rgb]{0.9,0.9,0.9}{
            \makecell[{{p{\linewidth}}}]{
                \texttt{\tiny{[GM$\rangle$P2]}}
                \texttt{Ditt meddelande har vidarebefordrats till den andra spelaren.} \\
\\ 
\\ 
\texttt{Du spelar för närvarande runda 3 av maximalt 12.} \\
\texttt{Ni har tillsammans ackumulerat 0 av 8 straff.} \\
\texttt{Den andra spelaren skickade detta meddelande:} \\
\\ 
\\ 
\texttt{"Låter bra! Jag börjar flytta C till (5,7). Då kör vi igång!"} \\
\\ 
\\ 
\texttt{Vad är ditt nästa kommando?} \\
            }
        }
    }
    & \\ \\

    \theutterance \stepcounter{utterance}  
    & & & \multicolumn{4}{p{0.6\linewidth}}{
        \cellcolor[rgb]{0.9,0.9,0.9}{
            \makecell[{{p{\linewidth}}}]{
                \texttt{\tiny{[GM$\langle$P2]}}
                \texttt{Låt mig först dubbelkolla spelplanen. Min L är på (2,2). Jag ska flytta den till (4,4). Jag måste kontrollera att (4,4) är en tom cell.} \\
\\ 
\texttt{Tittar på rutnätet, rad 4, kolumn 4: I rad 4 ser jag `║◌┌─┬┴┬─╢`. Kolumn 4 har symbolen `┴` {-} det är en linje, inte tom. Så (4,4) är inte ledigt.} \\
\\ 
\texttt{Jag behöver informera min partner om detta.} \\
\\ 
\texttt{SÄG: Vänta! Jag kollade min plan och (4,4) är inte en tom cell {-} där finns en linje (┴). Jag kan inte flytta L dit. Kan vi välja en annan position för L? (4,5) kanske? Eller har du något annat förslag?} \\
            }
        }
    }
     \\ \\

    \theutterance \stepcounter{utterance}  
    & & & \multicolumn{2}{p{0.3\linewidth}}{
        \cellcolor[rgb]{0.9,0.9,0.9}{
            \makecell[{{p{\linewidth}}}]{
                \texttt{\tiny{[GM$|$GM]}}
                \texttt{Låt mig först dubbelkolla spelplanen. Min L är på (2,2). Jag ska flytta den till (4,4). Jag måste kontrollera att (4,4) är en tom cell.} \\
\\ 
\texttt{Tittar på rutnätet, rad 4, kolumn 4: I rad 4 ser jag `║◌┌─┬┴┬─╢`. Kolumn 4 har symbolen `┴` {-} det är en linje, inte tom. Så (4,4) är inte ledigt.} \\
\\ 
\texttt{Jag behöver informera min partner om detta.} \\
\\ 
\texttt{SÄG: Vänta! Jag kollade min plan och (4,4) är inte en tom cell {-} där finns en linje (┴). Jag kan inte flytta L dit. Kan vi välja en annan position för L? (4,5) kanske? Eller har du något annat förslag?} \\
            }
        }
    }
    & & \\ \\

    \theutterance \stepcounter{utterance}  
    & & & \multicolumn{2}{p{0.3\linewidth}}{
        \cellcolor[rgb]{0.9,0.9,0.9}{
            \makecell[{{p{\linewidth}}}]{
                \texttt{\tiny{[GM$|$GM]}}
                \texttt{Invalid format: head is not empty: 'Låt mig först dubbelkolla spelplanen. Min L är på (2,2). Jag ska flytta den till (4,4). Jag måste kontrollera att (4,4) är en tom cell.} \\
\\ 
\texttt{Tittar på rutnätet, rad 4, kolumn 4: I rad 4 ser jag ║◌┌─┬┴┬─╢. Kolumn 4 har symbolen ┴ {-} det är en linje, inte tom. Så (4,4) är inte ledigt.} \\
\\ 
\texttt{Jag behöver informera min partner om detta.} \\
\\ 
\texttt{'} \\
            }
        }
    }
    & & \\ \\

    \theutterance \stepcounter{utterance}  
    & & \multicolumn{4}{p{0.6\linewidth}}{
        \cellcolor[rgb]{0.9,0.9,0.9}{
            \makecell[{{p{\linewidth}}}]{
                \texttt{\tiny{[GM$\rangle$P2]}}
                \texttt{Straff: Ditt meddelande får inte innehålla något före kommandot!} \\
\\ 
\texttt{Se till att ditt svar endast innehåller antingen `SÄG: <MEDDELANDE>` eller `FLYTTA: <OBJEKT>, (<X>, <Y>)`, och inget annat!} \\
\\ 
\texttt{Ni har tillsammans ackumulerat 1 av 8 straff.} \\
\texttt{Vänligen försök igen!} \\
            }
        }
    }
    & \\ \\

    \theutterance \stepcounter{utterance}  
    & & & \multicolumn{4}{p{0.6\linewidth}}{
        \cellcolor[rgb]{0.9,0.9,0.9}{
            \makecell[{{p{\linewidth}}}]{
                \texttt{\tiny{[GM$\langle$P2]}}
                \texttt{SÄG: Vänta! Jag kollade min plan och (4,4) är inte en tom cell {-} där finns en linje (┴). Jag kan inte flytta L dit. Kan vi välja en annan position för L? (4,5) kanske? Eller har du något annat förslag?} \\
            }
        }
    }
     \\ \\

    \theutterance \stepcounter{utterance}  
    & & & \multicolumn{2}{p{0.3\linewidth}}{
        \cellcolor[rgb]{0.9,0.9,0.9}{
            \makecell[{{p{\linewidth}}}]{
                \texttt{\tiny{[GM$|$GM]}}
                \texttt{SÄG: Vänta! Jag kollade min plan och (4,4) är inte en tom cell {-} där finns en linje (┴). Jag kan inte flytta L dit. Kan vi välja en annan position för L? (4,5) kanske? Eller har du något annat förslag?} \\
            }
        }
    }
    & & \\ \\

    \theutterance \stepcounter{utterance}  
    & & \multicolumn{4}{p{0.6\linewidth}}{
        \cellcolor[rgb]{0.9,0.9,0.9}{
            \makecell[{{p{\linewidth}}}]{
                \texttt{\tiny{[P1$\langle$GM]}}
                \texttt{Ditt meddelande har vidarebefordrats till den andra spelaren.} \\
\\ 
\\ 
\texttt{Du spelar för närvarande runda 3 av maximalt 12.} \\
\texttt{Ni har tillsammans ackumulerat 1 av 8 straff.} \\
\texttt{Den andra spelaren skickade detta meddelande:} \\
\\ 
\\ 
\texttt{"Vänta! Jag kollade min plan och (4,4) är inte en tom cell {-} där finns en linje (┴). Jag kan inte flytta L dit. Kan vi välja en annan position för L? (4,5) kanske? Eller har du något annat förslag?"} \\
\\ 
\\ 
\texttt{Vad är ditt nästa kommando?} \\
            }
        }
    }
    & \\ \\

    \theutterance \stepcounter{utterance}  
    & \multicolumn{4}{p{0.6\linewidth}}{
        \cellcolor[rgb]{0.9,0.9,0.9}{
            \makecell[{{p{\linewidth}}}]{
                \texttt{\tiny{[P1$\rangle$GM]}}
                \texttt{SÄG: Oj, bra att du kollade! Låt mig kolla min plan... (4,4) är också upptaget hos mig med en linje (┬). (4,5) är däremot tomt hos mig. Kör på (4,5) för L!} \\
            }
        }
    }
    & & \\ \\

    \theutterance \stepcounter{utterance}  
    & & & \multicolumn{2}{p{0.3\linewidth}}{
        \cellcolor[rgb]{0.9,0.9,0.9}{
            \makecell[{{p{\linewidth}}}]{
                \texttt{\tiny{[GM$|$GM]}}
                \texttt{SÄG: Oj, bra att du kollade! Låt mig kolla min plan... (4,4) är också upptaget hos mig med en linje (┬). (4,5) är däremot tomt hos mig. Kör på (4,5) för L!} \\
            }
        }
    }
    & & \\ \\

    \theutterance \stepcounter{utterance}  
    & & \multicolumn{4}{p{0.6\linewidth}}{
        \cellcolor[rgb]{0.9,0.9,0.9}{
            \makecell[{{p{\linewidth}}}]{
                \texttt{\tiny{[GM$\rangle$P2]}}
                \texttt{Ditt meddelande har vidarebefordrats till den andra spelaren.} \\
\\ 
\\ 
\texttt{Du spelar för närvarande runda 4 av maximalt 12.} \\
\texttt{Ni har tillsammans ackumulerat 1 av 8 straff.} \\
\texttt{Den andra spelaren skickade detta meddelande:} \\
\\ 
\\ 
\texttt{"Oj, bra att du kollade! Låt mig kolla min plan... (4,4) är också upptaget hos mig med en linje (┬). (4,5) är däremot tomt hos mig. Kör på (4,5) för L!"} \\
\\ 
\\ 
\texttt{Vad är ditt nästa kommando?} \\
            }
        }
    }
    & \\ \\

    \theutterance \stepcounter{utterance}  
    & & & \multicolumn{4}{p{0.6\linewidth}}{
        \cellcolor[rgb]{0.9,0.9,0.9}{
            \makecell[{{p{\linewidth}}}]{
                \texttt{\tiny{[GM$\langle$P2]}}
                \texttt{Låt mig kontrollera (4,5) på min plan. Rad 5: `╟─┤◌│◌│◌║`. Kolumn 4 på rad 5 är `◌` {-} den är tom! Jag flyttar L från (2,2) till (4,5).} \\
\\ 
\texttt{FLYTTA: L, (4,5)} \\
            }
        }
    }
     \\ \\

    \theutterance \stepcounter{utterance}  
    & & & \multicolumn{2}{p{0.3\linewidth}}{
        \cellcolor[rgb]{0.9,0.9,0.9}{
            \makecell[{{p{\linewidth}}}]{
                \texttt{\tiny{[GM$|$GM]}}
                \texttt{Låt mig kontrollera (4,5) på min plan. Rad 5: `╟─┤◌│◌│◌║`. Kolumn 4 på rad 5 är `◌` {-} den är tom! Jag flyttar L från (2,2) till (4,5).} \\
\\ 
\texttt{FLYTTA: L, (4,5)} \\
            }
        }
    }
    & & \\ \\

    \theutterance \stepcounter{utterance}  
    & & & \multicolumn{2}{p{0.3\linewidth}}{
        \cellcolor[rgb]{0.9,0.9,0.9}{
            \makecell[{{p{\linewidth}}}]{
                \texttt{\tiny{[GM$|$GM]}}
                \texttt{Invalid format: head is not empty: 'Låt mig kontrollera (4,5) på min plan. Rad 5: ╟─┤◌│◌│◌║. Kolumn 4 på rad 5 är ◌ {-} den är tom! Jag flyttar L från (2,2) till (4,5).} \\
\\ 
\texttt{'} \\
            }
        }
    }
    & & \\ \\

    \theutterance \stepcounter{utterance}  
    & & \multicolumn{4}{p{0.6\linewidth}}{
        \cellcolor[rgb]{0.9,0.9,0.9}{
            \makecell[{{p{\linewidth}}}]{
                \texttt{\tiny{[GM$\rangle$P2]}}
                \texttt{Straff: Ditt meddelande får inte innehålla något före kommandot!} \\
\\ 
\texttt{Se till att ditt svar endast innehåller antingen `SÄG: <MEDDELANDE>` eller `FLYTTA: <OBJEKT>, (<X>, <Y>)`, och inget annat!} \\
\\ 
\texttt{Ni har tillsammans ackumulerat 2 av 8 straff.} \\
\texttt{Vänligen försök igen!} \\
            }
        }
    }
    & \\ \\

    \theutterance \stepcounter{utterance}  
    & & & \multicolumn{4}{p{0.6\linewidth}}{
        \cellcolor[rgb]{0.9,0.9,0.9}{
            \makecell[{{p{\linewidth}}}]{
                \texttt{\tiny{[GM$\langle$P2]}}
                \texttt{FLYTTA: L, (4,5)} \\
            }
        }
    }
     \\ \\

    \theutterance \stepcounter{utterance}  
    & & & \multicolumn{2}{p{0.3\linewidth}}{
        \cellcolor[rgb]{0.9,0.9,0.9}{
            \makecell[{{p{\linewidth}}}]{
                \texttt{\tiny{[GM$|$GM]}}
                \texttt{FLYTTA: L, (4,5)} \\
            }
        }
    }
    & & \\ \\

    \theutterance \stepcounter{utterance}  
    & & & \multicolumn{2}{p{0.3\linewidth}}{
        \cellcolor[rgb]{0.9,0.9,0.9}{
            \makecell[{{p{\linewidth}}}]{
                \texttt{\tiny{[GM$|$GM]}}
                \texttt{Straff: (4,5) är inte tomt, utan innehåller '│'.} \\
\\ 
\texttt{Ni har tillsammans ackumulerat 3 av 8 straff.} \\
\texttt{Vänligen försök igen!} \\
            }
        }
    }
    & & \\ \\

    \theutterance \stepcounter{utterance}  
    & & \multicolumn{4}{p{0.6\linewidth}}{
        \cellcolor[rgb]{0.9,0.9,0.9}{
            \makecell[{{p{\linewidth}}}]{
                \texttt{\tiny{[GM$\rangle$P2]}}
                \texttt{Straff: (4,5) är inte tomt, utan innehåller '│'.} \\
\\ 
\texttt{Ni har tillsammans ackumulerat 3 av 8 straff.} \\
\texttt{Vänligen försök igen!} \\
            }
        }
    }
    & \\ \\

    \theutterance \stepcounter{utterance}  
    & & & \multicolumn{4}{p{0.6\linewidth}}{
        \cellcolor[rgb]{0.9,0.9,0.9}{
            \makecell[{{p{\linewidth}}}]{
                \texttt{\tiny{[GM$\langle$P2]}}
                \texttt{SÄG: (4,5) är inte tomt på min plan heller, där finns en linje (│). Låt oss hitta en position som är tom hos oss båda. Kan du föreslå en annan plats för L?} \\
            }
        }
    }
     \\ \\

    \theutterance \stepcounter{utterance}  
    & & & \multicolumn{2}{p{0.3\linewidth}}{
        \cellcolor[rgb]{0.9,0.9,0.9}{
            \makecell[{{p{\linewidth}}}]{
                \texttt{\tiny{[GM$|$GM]}}
                \texttt{SÄG: (4,5) är inte tomt på min plan heller, där finns en linje (│). Låt oss hitta en position som är tom hos oss båda. Kan du föreslå en annan plats för L?} \\
            }
        }
    }
    & & \\ \\

    \theutterance \stepcounter{utterance}  
    & & \multicolumn{4}{p{0.6\linewidth}}{
        \cellcolor[rgb]{0.9,0.9,0.9}{
            \makecell[{{p{\linewidth}}}]{
                \texttt{\tiny{[P1$\langle$GM]}}
                \texttt{Ditt meddelande har vidarebefordrats till den andra spelaren.} \\
\\ 
\\ 
\texttt{Du spelar för närvarande runda 4 av maximalt 12.} \\
\texttt{Ni har tillsammans ackumulerat 3 av 8 straff.} \\
\texttt{Den andra spelaren skickade detta meddelande:} \\
\\ 
\\ 
\texttt{"(4,5) är inte tomt på min plan heller, där finns en linje (│). Låt oss hitta en position som är tom hos oss båda. Kan du föreslå en annan plats för L?"} \\
\\ 
\\ 
\texttt{Vad är ditt nästa kommando?} \\
            }
        }
    }
    & \\ \\

    \theutterance \stepcounter{utterance}  
    & \multicolumn{4}{p{0.6\linewidth}}{
        \cellcolor[rgb]{0.9,0.9,0.9}{
            \makecell[{{p{\linewidth}}}]{
                \texttt{\tiny{[P1$\rangle$GM]}}
                \texttt{SÄG: Låt oss tänka om. Min L är vid (6,6) och din vid (2,2). Vad sägs om (5,5)? Jag kollar min plan... (5,5) är tomt hos mig (◌). Är det tomt hos dig?} \\
            }
        }
    }
    & & \\ \\

    \theutterance \stepcounter{utterance}  
    & & & \multicolumn{2}{p{0.3\linewidth}}{
        \cellcolor[rgb]{0.9,0.9,0.9}{
            \makecell[{{p{\linewidth}}}]{
                \texttt{\tiny{[GM$|$GM]}}
                \texttt{SÄG: Låt oss tänka om. Min L är vid (6,6) och din vid (2,2). Vad sägs om (5,5)? Jag kollar min plan... (5,5) är tomt hos mig (◌). Är det tomt hos dig?} \\
            }
        }
    }
    & & \\ \\

    \theutterance \stepcounter{utterance}  
    & & \multicolumn{4}{p{0.6\linewidth}}{
        \cellcolor[rgb]{0.9,0.9,0.9}{
            \makecell[{{p{\linewidth}}}]{
                \texttt{\tiny{[GM$\rangle$P2]}}
                \texttt{Ditt meddelande har vidarebefordrats till den andra spelaren.} \\
\\ 
\\ 
\texttt{Du spelar för närvarande runda 5 av maximalt 12.} \\
\texttt{Ni har tillsammans ackumulerat 3 av 8 straff.} \\
\texttt{Den andra spelaren skickade detta meddelande:} \\
\\ 
\\ 
\texttt{"Låt oss tänka om. Min L är vid (6,6) och din vid (2,2). Vad sägs om (5,5)? Jag kollar min plan... (5,5) är tomt hos mig (◌). Är det tomt hos dig?"} \\
\\ 
\\ 
\texttt{Vad är ditt nästa kommando?} \\
            }
        }
    }
    & \\ \\

    \theutterance \stepcounter{utterance}  
    & & & \multicolumn{4}{p{0.6\linewidth}}{
        \cellcolor[rgb]{0.9,0.9,0.9}{
            \makecell[{{p{\linewidth}}}]{
                \texttt{\tiny{[GM$\langle$P2]}}
                \texttt{Låt mig kolla (5,5) på min plan. Rad 5: `╟─┤◌│◌│◌║`. Kolumn 5 på rad 5 är `◌` {-} den är tom!} \\
\\ 
\texttt{SÄG: Ja, (5,5) är tomt hos mig! Jag flyttar L dit nu.} \\
            }
        }
    }
     \\ \\

    \theutterance \stepcounter{utterance}  
    & & & \multicolumn{2}{p{0.3\linewidth}}{
        \cellcolor[rgb]{0.9,0.9,0.9}{
            \makecell[{{p{\linewidth}}}]{
                \texttt{\tiny{[GM$|$GM]}}
                \texttt{Låt mig kolla (5,5) på min plan. Rad 5: `╟─┤◌│◌│◌║`. Kolumn 5 på rad 5 är `◌` {-} den är tom!} \\
\\ 
\texttt{SÄG: Ja, (5,5) är tomt hos mig! Jag flyttar L dit nu.} \\
            }
        }
    }
    & & \\ \\

    \theutterance \stepcounter{utterance}  
    & & & \multicolumn{2}{p{0.3\linewidth}}{
        \cellcolor[rgb]{0.9,0.9,0.9}{
            \makecell[{{p{\linewidth}}}]{
                \texttt{\tiny{[GM$|$GM]}}
                \texttt{Invalid format: head is not empty: 'Låt mig kolla (5,5) på min plan. Rad 5: ╟─┤◌│◌│◌║. Kolumn 5 på rad 5 är ◌ {-} den är tom!} \\
\\ 
\texttt{'} \\
            }
        }
    }
    & & \\ \\

    \theutterance \stepcounter{utterance}  
    & & \multicolumn{4}{p{0.6\linewidth}}{
        \cellcolor[rgb]{0.9,0.9,0.9}{
            \makecell[{{p{\linewidth}}}]{
                \texttt{\tiny{[GM$\rangle$P2]}}
                \texttt{Straff: Ditt meddelande får inte innehålla något före kommandot!} \\
\\ 
\texttt{Se till att ditt svar endast innehåller antingen `SÄG: <MEDDELANDE>` eller `FLYTTA: <OBJEKT>, (<X>, <Y>)`, och inget annat!} \\
\\ 
\texttt{Ni har tillsammans ackumulerat 4 av 8 straff.} \\
\texttt{Vänligen försök igen!} \\
            }
        }
    }
    & \\ \\

    \theutterance \stepcounter{utterance}  
    & & & \multicolumn{4}{p{0.6\linewidth}}{
        \cellcolor[rgb]{0.9,0.9,0.9}{
            \makecell[{{p{\linewidth}}}]{
                \texttt{\tiny{[GM$\langle$P2]}}
                \texttt{SÄG: Ja, (5,5) är tomt hos mig! Jag flyttar L dit nu.} \\
            }
        }
    }
     \\ \\

    \theutterance \stepcounter{utterance}  
    & & & \multicolumn{2}{p{0.3\linewidth}}{
        \cellcolor[rgb]{0.9,0.9,0.9}{
            \makecell[{{p{\linewidth}}}]{
                \texttt{\tiny{[GM$|$GM]}}
                \texttt{SÄG: Ja, (5,5) är tomt hos mig! Jag flyttar L dit nu.} \\
            }
        }
    }
    & & \\ \\

    \theutterance \stepcounter{utterance}  
    & & \multicolumn{4}{p{0.6\linewidth}}{
        \cellcolor[rgb]{0.9,0.9,0.9}{
            \makecell[{{p{\linewidth}}}]{
                \texttt{\tiny{[P1$\langle$GM]}}
                \texttt{Ditt meddelande har vidarebefordrats till den andra spelaren.} \\
\\ 
\\ 
\texttt{Du spelar för närvarande runda 5 av maximalt 12.} \\
\texttt{Ni har tillsammans ackumulerat 4 av 8 straff.} \\
\texttt{Den andra spelaren skickade detta meddelande:} \\
\\ 
\\ 
\texttt{"Ja, (5,5) är tomt hos mig! Jag flyttar L dit nu."} \\
\\ 
\\ 
\texttt{Vad är ditt nästa kommando?} \\
            }
        }
    }
    & \\ \\

    \theutterance \stepcounter{utterance}  
    & \multicolumn{4}{p{0.6\linewidth}}{
        \cellcolor[rgb]{0.9,0.9,0.9}{
            \makecell[{{p{\linewidth}}}]{
                \texttt{\tiny{[P1$\rangle$GM]}}
                \texttt{FLYTTA: L, (5,5)} \\
            }
        }
    }
    & & \\ \\

    \theutterance \stepcounter{utterance}  
    & & & \multicolumn{2}{p{0.3\linewidth}}{
        \cellcolor[rgb]{0.9,0.9,0.9}{
            \makecell[{{p{\linewidth}}}]{
                \texttt{\tiny{[GM$|$GM]}}
                \texttt{FLYTTA: L, (5,5)} \\
            }
        }
    }
    & & \\ \\

    \theutterance \stepcounter{utterance}  
    & & & \multicolumn{2}{p{0.3\linewidth}}{
        \cellcolor[rgb]{0.9,0.9,0.9}{
            \makecell[{{p{\linewidth}}}]{
                \texttt{\tiny{[GM$|$GM]}}
                \texttt{L moved to (5, 5)} \\
            }
        }
    }
    & & \\ \\

    \theutterance \stepcounter{utterance}  
    & & \multicolumn{4}{p{0.6\linewidth}}{
        \cellcolor[rgb]{0.9,0.9,0.9}{
            \makecell[{{p{\linewidth}}}]{
                \texttt{\tiny{[GM$\rangle$P2]}}
                \texttt{Ditt meddelande har vidarebefordrats till den andra spelaren.} \\
\\ 
\\ 
\texttt{Du spelar för närvarande runda 6 av maximalt 12.} \\
\texttt{Ni har tillsammans ackumulerat 4 av 8 straff.} \\
\texttt{Den andra spelaren flyttade ett objekt på sitt rutnät.} \\
\\ 
\\ 
\texttt{Vad är ditt nästa kommando?} \\
            }
        }
    }
    & \\ \\

    \theutterance \stepcounter{utterance}  
    & & & \multicolumn{4}{p{0.6\linewidth}}{
        \cellcolor[rgb]{0.9,0.9,0.9}{
            \makecell[{{p{\linewidth}}}]{
                \texttt{\tiny{[GM$\langle$P2]}}
                \texttt{FLYTTA: L, (5,5)} \\
            }
        }
    }
     \\ \\

    \theutterance \stepcounter{utterance}  
    & & & \multicolumn{2}{p{0.3\linewidth}}{
        \cellcolor[rgb]{0.9,0.9,0.9}{
            \makecell[{{p{\linewidth}}}]{
                \texttt{\tiny{[GM$|$GM]}}
                \texttt{FLYTTA: L, (5,5)} \\
            }
        }
    }
    & & \\ \\

    \theutterance \stepcounter{utterance}  
    & & & \multicolumn{2}{p{0.3\linewidth}}{
        \cellcolor[rgb]{0.9,0.9,0.9}{
            \makecell[{{p{\linewidth}}}]{
                \texttt{\tiny{[GM$|$GM]}}
                \texttt{L moved to (5, 5)} \\
            }
        }
    }
    & & \\ \\

    \theutterance \stepcounter{utterance}  
    & & \multicolumn{4}{p{0.6\linewidth}}{
        \cellcolor[rgb]{0.9,0.9,0.9}{
            \makecell[{{p{\linewidth}}}]{
                \texttt{\tiny{[P1$\langle$GM]}}
                \texttt{Flyttade 'L' till (5,5) framgångsrikt. Ditt uppdaterade rutnät ser ut så här:} \\
\texttt{```} \\
\texttt{ 1234567} \\
\texttt{╔══╤═╤══╗} \\
\texttt{║◌◌│◌│◌◌║ 1} \\
\texttt{║◌◌└─┼──╢ 2} \\
\texttt{║◌◌◌◌│P◌║ 3} \\
\texttt{║◌┌─┬┴┬─╢ 4} \\
\texttt{╟─┤◌│L│◌║ 5} \\
\texttt{║◌└┬┴┬┘◌║ 6} \\
\texttt{║◌◌│C│◌◌║ 7} \\
\texttt{╚══╧═╧══╝} \\
\texttt{```} \\
\texttt{Du spelar för närvarande runda 6 av maximalt 12.} \\
\texttt{Ni har tillsammans ackumulerat 4 av 8 straff.} \\
\texttt{Den andra spelaren flyttade ett objekt på sitt rutnät.} \\
\\ 
\\ 
\texttt{Vad är ditt nästa kommando?} \\
            }
        }
    }
    & \\ \\

    \theutterance \stepcounter{utterance}  
    & \multicolumn{4}{p{0.6\linewidth}}{
        \cellcolor[rgb]{0.9,0.9,0.9}{
            \makecell[{{p{\linewidth}}}]{
                \texttt{\tiny{[P1$\rangle$GM]}}
                \texttt{SÄG: Jag har flyttat L till (5,5). Nu flyttar jag C till (5,7). Gör du samma sak?} \\
            }
        }
    }
    & & \\ \\

    \theutterance \stepcounter{utterance}  
    & & & \multicolumn{2}{p{0.3\linewidth}}{
        \cellcolor[rgb]{0.9,0.9,0.9}{
            \makecell[{{p{\linewidth}}}]{
                \texttt{\tiny{[GM$|$GM]}}
                \texttt{SÄG: Jag har flyttat L till (5,5). Nu flyttar jag C till (5,7). Gör du samma sak?} \\
            }
        }
    }
    & & \\ \\

    \theutterance \stepcounter{utterance}  
    & & \multicolumn{4}{p{0.6\linewidth}}{
        \cellcolor[rgb]{0.9,0.9,0.9}{
            \makecell[{{p{\linewidth}}}]{
                \texttt{\tiny{[GM$\rangle$P2]}}
                \texttt{Flyttade 'L' till (5,5) framgångsrikt. Ditt uppdaterade rutnät ser ut så här:} \\
\texttt{```} \\
\texttt{ 1234567} \\
\texttt{╔══╤═╤══╗} \\
\texttt{║◌P│◌│◌◌║ 1} \\
\texttt{║◌◌└─┼──╢ 2} \\
\texttt{║◌◌◌◌│◌◌║ 3} \\
\texttt{║◌┌─┬┴┬─╢ 4} \\
\texttt{╟─┤◌│L│◌║ 5} \\
\texttt{║◌└┬┴┬┘◌║ 6} \\
\texttt{║◌◌│◌│◌C║ 7} \\
\texttt{╚══╧═╧══╝} \\
\texttt{```} \\
\texttt{Du spelar för närvarande runda 7 av maximalt 12.} \\
\texttt{Ni har tillsammans ackumulerat 4 av 8 straff.} \\
\texttt{Den andra spelaren skickade detta meddelande:} \\
\\ 
\\ 
\texttt{"Jag har flyttat L till (5,5). Nu flyttar jag C till (5,7). Gör du samma sak?"} \\
\\ 
\\ 
\texttt{Vad är ditt nästa kommando?} \\
            }
        }
    }
    & \\ \\

    \theutterance \stepcounter{utterance}  
    & & & \multicolumn{4}{p{0.6\linewidth}}{
        \cellcolor[rgb]{0.9,0.9,0.9}{
            \makecell[{{p{\linewidth}}}]{
                \texttt{\tiny{[GM$\langle$P2]}}
                \texttt{SÄG: Jag har också flyttat L till (5,5). För C {-} din är på (4,7) och min på (7,7). (5,7) är en bra kompromiss. Låt mig bara kolla att det är tomt hos mig först.} \\
            }
        }
    }
     \\ \\

    \theutterance \stepcounter{utterance}  
    & & & \multicolumn{2}{p{0.3\linewidth}}{
        \cellcolor[rgb]{0.9,0.9,0.9}{
            \makecell[{{p{\linewidth}}}]{
                \texttt{\tiny{[GM$|$GM]}}
                \texttt{SÄG: Jag har också flyttat L till (5,5). För C {-} din är på (4,7) och min på (7,7). (5,7) är en bra kompromiss. Låt mig bara kolla att det är tomt hos mig först.} \\
            }
        }
    }
    & & \\ \\

    \theutterance \stepcounter{utterance}  
    & & \multicolumn{4}{p{0.6\linewidth}}{
        \cellcolor[rgb]{0.9,0.9,0.9}{
            \makecell[{{p{\linewidth}}}]{
                \texttt{\tiny{[P1$\langle$GM]}}
                \texttt{Ditt meddelande har vidarebefordrats till den andra spelaren.} \\
\\ 
\\ 
\texttt{Du spelar för närvarande runda 7 av maximalt 12.} \\
\texttt{Ni har tillsammans ackumulerat 4 av 8 straff.} \\
\texttt{Den andra spelaren skickade detta meddelande:} \\
\\ 
\\ 
\texttt{"Jag har också flyttat L till (5,5). För C {-} din är på (4,7) och min på (7,7). (5,7) är en bra kompromiss. Låt mig bara kolla att det är tomt hos mig först."} \\
\\ 
\\ 
\texttt{Vad är ditt nästa kommando?} \\
            }
        }
    }
    & \\ \\

    \theutterance \stepcounter{utterance}  
    & \multicolumn{4}{p{0.6\linewidth}}{
        \cellcolor[rgb]{0.9,0.9,0.9}{
            \makecell[{{p{\linewidth}}}]{
                \texttt{\tiny{[P1$\rangle$GM]}}
                \texttt{SÄG: Bra! (5,7) är tomt hos mig {-} där finns bara ◌. Jag flyttar C dit nu!} \\
            }
        }
    }
    & & \\ \\

    \theutterance \stepcounter{utterance}  
    & & & \multicolumn{2}{p{0.3\linewidth}}{
        \cellcolor[rgb]{0.9,0.9,0.9}{
            \makecell[{{p{\linewidth}}}]{
                \texttt{\tiny{[GM$|$GM]}}
                \texttt{SÄG: Bra! (5,7) är tomt hos mig {-} där finns bara ◌. Jag flyttar C dit nu!} \\
            }
        }
    }
    & & \\ \\

    \theutterance \stepcounter{utterance}  
    & & \multicolumn{4}{p{0.6\linewidth}}{
        \cellcolor[rgb]{0.9,0.9,0.9}{
            \makecell[{{p{\linewidth}}}]{
                \texttt{\tiny{[GM$\rangle$P2]}}
                \texttt{Ditt meddelande har vidarebefordrats till den andra spelaren.} \\
\\ 
\\ 
\texttt{Du spelar för närvarande runda 8 av maximalt 12.} \\
\texttt{Ni har tillsammans ackumulerat 4 av 8 straff.} \\
\texttt{Den andra spelaren skickade detta meddelande:} \\
\\ 
\\ 
\texttt{"Bra! (5,7) är tomt hos mig {-} där finns bara ◌. Jag flyttar C dit nu!"} \\
\\ 
\\ 
\texttt{Vad är ditt nästa kommando?} \\
            }
        }
    }
    & \\ \\

    \theutterance \stepcounter{utterance}  
    & & & \multicolumn{4}{p{0.6\linewidth}}{
        \cellcolor[rgb]{0.9,0.9,0.9}{
            \makecell[{{p{\linewidth}}}]{
                \texttt{\tiny{[GM$\langle$P2]}}
                \texttt{Låt mig kolla (5,7) på min plan. Rad 7: `║◌◌│◌│◌C║`. Kolumn 5 på rad 7 är `◌` {-} den är tom!} \\
\\ 
\texttt{FLYTTA: C, (5,7)} \\
            }
        }
    }
     \\ \\

    \theutterance \stepcounter{utterance}  
    & & & \multicolumn{2}{p{0.3\linewidth}}{
        \cellcolor[rgb]{0.9,0.9,0.9}{
            \makecell[{{p{\linewidth}}}]{
                \texttt{\tiny{[GM$|$GM]}}
                \texttt{Låt mig kolla (5,7) på min plan. Rad 7: `║◌◌│◌│◌C║`. Kolumn 5 på rad 7 är `◌` {-} den är tom!} \\
\\ 
\texttt{FLYTTA: C, (5,7)} \\
            }
        }
    }
    & & \\ \\

    \theutterance \stepcounter{utterance}  
    & & & \multicolumn{2}{p{0.3\linewidth}}{
        \cellcolor[rgb]{0.9,0.9,0.9}{
            \makecell[{{p{\linewidth}}}]{
                \texttt{\tiny{[GM$|$GM]}}
                \texttt{Invalid format: head is not empty: 'Låt mig kolla (5,7) på min plan. Rad 7: ║◌◌│◌│◌C║. Kolumn 5 på rad 7 är ◌ {-} den är tom!} \\
\\ 
\texttt{'} \\
            }
        }
    }
    & & \\ \\

    \theutterance \stepcounter{utterance}  
    & & \multicolumn{4}{p{0.6\linewidth}}{
        \cellcolor[rgb]{0.9,0.9,0.9}{
            \makecell[{{p{\linewidth}}}]{
                \texttt{\tiny{[GM$\rangle$P2]}}
                \texttt{Straff: Ditt meddelande får inte innehålla något före kommandot!} \\
\\ 
\texttt{Se till att ditt svar endast innehåller antingen `SÄG: <MEDDELANDE>` eller `FLYTTA: <OBJEKT>, (<X>, <Y>)`, och inget annat!} \\
\\ 
\texttt{Ni har tillsammans ackumulerat 5 av 8 straff.} \\
\texttt{Vänligen försök igen!} \\
            }
        }
    }
    & \\ \\

    \theutterance \stepcounter{utterance}  
    & & & \multicolumn{4}{p{0.6\linewidth}}{
        \cellcolor[rgb]{0.9,0.9,0.9}{
            \makecell[{{p{\linewidth}}}]{
                \texttt{\tiny{[GM$\langle$P2]}}
                \texttt{FLYTTA: C, (5,7)} \\
            }
        }
    }
     \\ \\

    \theutterance \stepcounter{utterance}  
    & & & \multicolumn{2}{p{0.3\linewidth}}{
        \cellcolor[rgb]{0.9,0.9,0.9}{
            \makecell[{{p{\linewidth}}}]{
                \texttt{\tiny{[GM$|$GM]}}
                \texttt{FLYTTA: C, (5,7)} \\
            }
        }
    }
    & & \\ \\

    \theutterance \stepcounter{utterance}  
    & & & \multicolumn{2}{p{0.3\linewidth}}{
        \cellcolor[rgb]{0.9,0.9,0.9}{
            \makecell[{{p{\linewidth}}}]{
                \texttt{\tiny{[GM$|$GM]}}
                \texttt{Straff: (5,7) är inte tomt, utan innehåller '│'.} \\
\\ 
\texttt{Ni har tillsammans ackumulerat 6 av 8 straff.} \\
\texttt{Vänligen försök igen!} \\
            }
        }
    }
    & & \\ \\

    \theutterance \stepcounter{utterance}  
    & & \multicolumn{4}{p{0.6\linewidth}}{
        \cellcolor[rgb]{0.9,0.9,0.9}{
            \makecell[{{p{\linewidth}}}]{
                \texttt{\tiny{[GM$\rangle$P2]}}
                \texttt{Straff: (5,7) är inte tomt, utan innehåller '│'.} \\
\\ 
\texttt{Ni har tillsammans ackumulerat 6 av 8 straff.} \\
\texttt{Vänligen försök igen!} \\
            }
        }
    }
    & \\ \\

    \theutterance \stepcounter{utterance}  
    & & & \multicolumn{4}{p{0.6\linewidth}}{
        \cellcolor[rgb]{0.9,0.9,0.9}{
            \makecell[{{p{\linewidth}}}]{
                \texttt{\tiny{[GM$\langle$P2]}}
                \texttt{SÄG: (5,7) är inte tomt hos mig, där finns en linje (│). Vi måste hitta en annan plats för C. Vad sägs om (6,7)? Låt mig kolla {-} nej, min C är redan på (7,7). Kanske (5,6) eller (4,6)? Kan du kolla din plan?} \\
            }
        }
    }
     \\ \\

    \theutterance \stepcounter{utterance}  
    & & & \multicolumn{2}{p{0.3\linewidth}}{
        \cellcolor[rgb]{0.9,0.9,0.9}{
            \makecell[{{p{\linewidth}}}]{
                \texttt{\tiny{[GM$|$GM]}}
                \texttt{SÄG: (5,7) är inte tomt hos mig, där finns en linje (│). Vi måste hitta en annan plats för C. Vad sägs om (6,7)? Låt mig kolla {-} nej, min C är redan på (7,7). Kanske (5,6) eller (4,6)? Kan du kolla din plan?} \\
            }
        }
    }
    & & \\ \\

    \theutterance \stepcounter{utterance}  
    & & \multicolumn{4}{p{0.6\linewidth}}{
        \cellcolor[rgb]{0.9,0.9,0.9}{
            \makecell[{{p{\linewidth}}}]{
                \texttt{\tiny{[P1$\langle$GM]}}
                \texttt{Ditt meddelande har vidarebefordrats till den andra spelaren.} \\
\\ 
\\ 
\texttt{Du spelar för närvarande runda 8 av maximalt 12.} \\
\texttt{Ni har tillsammans ackumulerat 6 av 8 straff.} \\
\texttt{Den andra spelaren skickade detta meddelande:} \\
\\ 
\\ 
\texttt{"(5,7) är inte tomt hos mig, där finns en linje (│). Vi måste hitta en annan plats för C. Vad sägs om (6,7)? Låt mig kolla {-} nej, min C är redan på (7,7). Kanske (5,6) eller (4,6)? Kan du kolla din plan?"} \\
\\ 
\\ 
\texttt{Vad är ditt nästa kommando?} \\
            }
        }
    }
    & \\ \\

    \theutterance \stepcounter{utterance}  
    & \multicolumn{4}{p{0.6\linewidth}}{
        \cellcolor[rgb]{0.9,0.9,0.9}{
            \makecell[{{p{\linewidth}}}]{
                \texttt{\tiny{[P1$\rangle$GM]}}
                \texttt{SÄG: Jag kollar min plan... (5,6) har en linje (┴) så den går inte. (4,6) däremot är tomt (◌) hos mig! Är (4,6) tomt hos dig?} \\
            }
        }
    }
    & & \\ \\

    \theutterance \stepcounter{utterance}  
    & & & \multicolumn{2}{p{0.3\linewidth}}{
        \cellcolor[rgb]{0.9,0.9,0.9}{
            \makecell[{{p{\linewidth}}}]{
                \texttt{\tiny{[GM$|$GM]}}
                \texttt{SÄG: Jag kollar min plan... (5,6) har en linje (┴) så den går inte. (4,6) däremot är tomt (◌) hos mig! Är (4,6) tomt hos dig?} \\
            }
        }
    }
    & & \\ \\

    \theutterance \stepcounter{utterance}  
    & & \multicolumn{4}{p{0.6\linewidth}}{
        \cellcolor[rgb]{0.9,0.9,0.9}{
            \makecell[{{p{\linewidth}}}]{
                \texttt{\tiny{[GM$\rangle$P2]}}
                \texttt{Ditt meddelande har vidarebefordrats till den andra spelaren.} \\
\\ 
\\ 
\texttt{Du spelar för närvarande runda 9 av maximalt 12.} \\
\texttt{Ni har tillsammans ackumulerat 6 av 8 straff.} \\
\texttt{Den andra spelaren skickade detta meddelande:} \\
\\ 
\\ 
\texttt{"Jag kollar min plan... (5,6) har en linje (┴) så den går inte. (4,6) däremot är tomt (◌) hos mig! Är (4,6) tomt hos dig?"} \\
\\ 
\\ 
\texttt{Vad är ditt nästa kommando?} \\
            }
        }
    }
    & \\ \\

    \theutterance \stepcounter{utterance}  
    & & & \multicolumn{4}{p{0.6\linewidth}}{
        \cellcolor[rgb]{0.9,0.9,0.9}{
            \makecell[{{p{\linewidth}}}]{
                \texttt{\tiny{[GM$\langle$P2]}}
                \texttt{Låt mig kolla (4,6) på min plan. Rad 6: `║◌└┬┴┬┘◌║`. Kolumn 4 på rad 6 är `┴` {-} en linje, inte tom.} \\
\\ 
\texttt{SÄG: (4,6) har en linje (┴) hos mig. Inte tomt. Låt oss testa (3,6) eller (6,6)? Hur ser de ut hos dig?} \\
            }
        }
    }
     \\ \\

    \theutterance \stepcounter{utterance}  
    & & & \multicolumn{2}{p{0.3\linewidth}}{
        \cellcolor[rgb]{0.9,0.9,0.9}{
            \makecell[{{p{\linewidth}}}]{
                \texttt{\tiny{[GM$|$GM]}}
                \texttt{Låt mig kolla (4,6) på min plan. Rad 6: `║◌└┬┴┬┘◌║`. Kolumn 4 på rad 6 är `┴` {-} en linje, inte tom.} \\
\\ 
\texttt{SÄG: (4,6) har en linje (┴) hos mig. Inte tomt. Låt oss testa (3,6) eller (6,6)? Hur ser de ut hos dig?} \\
            }
        }
    }
    & & \\ \\

    \theutterance \stepcounter{utterance}  
    & & & \multicolumn{2}{p{0.3\linewidth}}{
        \cellcolor[rgb]{0.9,0.9,0.9}{
            \makecell[{{p{\linewidth}}}]{
                \texttt{\tiny{[GM$|$GM]}}
                \texttt{Invalid format: head is not empty: 'Låt mig kolla (4,6) på min plan. Rad 6: ║◌└┬┴┬┘◌║. Kolumn 4 på rad 6 är ┴ {-} en linje, inte tom.} \\
\\ 
\texttt{'} \\
            }
        }
    }
    & & \\ \\

    \theutterance \stepcounter{utterance}  
    & & \multicolumn{4}{p{0.6\linewidth}}{
        \cellcolor[rgb]{0.9,0.9,0.9}{
            \makecell[{{p{\linewidth}}}]{
                \texttt{\tiny{[GM$\rangle$P2]}}
                \texttt{Straff: Ditt meddelande får inte innehålla något före kommandot!} \\
\\ 
\texttt{Se till att ditt svar endast innehåller antingen `SÄG: <MEDDELANDE>` eller `FLYTTA: <OBJEKT>, (<X>, <Y>)`, och inget annat!} \\
\\ 
\texttt{Ni har tillsammans ackumulerat 7 av 8 straff.} \\
\texttt{Vänligen försök igen!} \\
            }
        }
    }
    & \\ \\

    \theutterance \stepcounter{utterance}  
    & & & \multicolumn{4}{p{0.6\linewidth}}{
        \cellcolor[rgb]{0.9,0.9,0.9}{
            \makecell[{{p{\linewidth}}}]{
                \texttt{\tiny{[GM$\langle$P2]}}
                \texttt{SÄG: (4,6) har en linje (┴) hos mig. Inte tomt. Låt oss testa (3,6) eller (6,6)? Hur ser de ut hos dig?} \\
            }
        }
    }
     \\ \\

    \theutterance \stepcounter{utterance}  
    & & & \multicolumn{2}{p{0.3\linewidth}}{
        \cellcolor[rgb]{0.9,0.9,0.9}{
            \makecell[{{p{\linewidth}}}]{
                \texttt{\tiny{[GM$|$GM]}}
                \texttt{SÄG: (4,6) har en linje (┴) hos mig. Inte tomt. Låt oss testa (3,6) eller (6,6)? Hur ser de ut hos dig?} \\
            }
        }
    }
    & & \\ \\

    \theutterance \stepcounter{utterance}  
    & & \multicolumn{4}{p{0.6\linewidth}}{
        \cellcolor[rgb]{0.9,0.9,0.9}{
            \makecell[{{p{\linewidth}}}]{
                \texttt{\tiny{[P1$\langle$GM]}}
                \texttt{Ditt meddelande har vidarebefordrats till den andra spelaren.} \\
\\ 
\\ 
\texttt{Du spelar för närvarande runda 9 av maximalt 12.} \\
\texttt{Ni har tillsammans ackumulerat 7 av 8 straff.} \\
\texttt{Den andra spelaren skickade detta meddelande:} \\
\\ 
\\ 
\texttt{"(4,6) har en linje (┴) hos mig. Inte tomt. Låt oss testa (3,6) eller (6,6)? Hur ser de ut hos dig?"} \\
\\ 
\\ 
\texttt{Vad är ditt nästa kommando?} \\
            }
        }
    }
    & \\ \\

    \theutterance \stepcounter{utterance}  
    & \multicolumn{4}{p{0.6\linewidth}}{
        \cellcolor[rgb]{0.9,0.9,0.9}{
            \makecell[{{p{\linewidth}}}]{
                \texttt{\tiny{[P1$\rangle$GM]}}
                \texttt{SÄG: (3,6) är tomt (◌) hos mig! (6,6) är också tomt nu eftersom jag flyttade L därifrån. Vilken föredrar du?} \\
            }
        }
    }
    & & \\ \\

    \theutterance \stepcounter{utterance}  
    & & & \multicolumn{2}{p{0.3\linewidth}}{
        \cellcolor[rgb]{0.9,0.9,0.9}{
            \makecell[{{p{\linewidth}}}]{
                \texttt{\tiny{[GM$|$GM]}}
                \texttt{SÄG: (3,6) är tomt (◌) hos mig! (6,6) är också tomt nu eftersom jag flyttade L därifrån. Vilken föredrar du?} \\
            }
        }
    }
    & & \\ \\

    \theutterance \stepcounter{utterance}  
    & & \multicolumn{4}{p{0.6\linewidth}}{
        \cellcolor[rgb]{0.9,0.9,0.9}{
            \makecell[{{p{\linewidth}}}]{
                \texttt{\tiny{[GM$\rangle$P2]}}
                \texttt{Ditt meddelande har vidarebefordrats till den andra spelaren.} \\
\\ 
\\ 
\texttt{Du spelar för närvarande runda 10 av maximalt 12.} \\
\texttt{Ni har tillsammans ackumulerat 7 av 8 straff.} \\
\texttt{Den andra spelaren skickade detta meddelande:} \\
\\ 
\\ 
\texttt{"(3,6) är tomt (◌) hos mig! (6,6) är också tomt nu eftersom jag flyttade L därifrån. Vilken föredrar du?"} \\
\\ 
\\ 
\texttt{Vad är ditt nästa kommando?} \\
            }
        }
    }
    & \\ \\

    \theutterance \stepcounter{utterance}  
    & & & \multicolumn{4}{p{0.6\linewidth}}{
        \cellcolor[rgb]{0.9,0.9,0.9}{
            \makecell[{{p{\linewidth}}}]{
                \texttt{\tiny{[GM$\langle$P2]}}
                \texttt{Låt mig kolla båda. Rad 6: `║◌└┬┴┬┘◌║`. Kolumn 3 på rad 6 är `└` {-} en linje. Kolumn 6 på rad 6 är `┘` {-} också en linje. Ingen av dem är tomma!} \\
\\ 
\texttt{SÄG: Både (3,6) och (6,6) har linjer hos mig (└ och ┘). Vad sägs om (3,7) eller (6,5)? Kan du kolla dem?} \\
            }
        }
    }
     \\ \\

    \theutterance \stepcounter{utterance}  
    & & & \multicolumn{2}{p{0.3\linewidth}}{
        \cellcolor[rgb]{0.9,0.9,0.9}{
            \makecell[{{p{\linewidth}}}]{
                \texttt{\tiny{[GM$|$GM]}}
                \texttt{Låt mig kolla båda. Rad 6: `║◌└┬┴┬┘◌║`. Kolumn 3 på rad 6 är `└` {-} en linje. Kolumn 6 på rad 6 är `┘` {-} också en linje. Ingen av dem är tomma!} \\
\\ 
\texttt{SÄG: Både (3,6) och (6,6) har linjer hos mig (└ och ┘). Vad sägs om (3,7) eller (6,5)? Kan du kolla dem?} \\
            }
        }
    }
    & & \\ \\

    \theutterance \stepcounter{utterance}  
    & & & \multicolumn{2}{p{0.3\linewidth}}{
        \cellcolor[rgb]{0.9,0.9,0.9}{
            \makecell[{{p{\linewidth}}}]{
                \texttt{\tiny{[GM$|$GM]}}
                \texttt{Invalid format: head is not empty: 'Låt mig kolla båda. Rad 6: ║◌└┬┴┬┘◌║. Kolumn 3 på rad 6 är └ {-} en linje. Kolumn 6 på rad 6 är ┘ {-} också en linje. Ingen av dem är tomma!} \\
\\ 
\texttt{'} \\
            }
        }
    }
    & & \\ \\

    \theutterance \stepcounter{utterance}  
    & & \multicolumn{4}{p{0.6\linewidth}}{
        \cellcolor[rgb]{0.9,0.9,0.9}{
            \makecell[{{p{\linewidth}}}]{
                \texttt{\tiny{[GM$\rangle$P2]}}
                \texttt{Straff: Ditt meddelande får inte innehålla något före kommandot!} \\
\\ 
\texttt{Se till att ditt svar endast innehåller antingen `SÄG: <MEDDELANDE>` eller `FLYTTA: <OBJEKT>, (<X>, <Y>)`, och inget annat!} \\
\\ 
\texttt{Ni har tillsammans ackumulerat 8 av 8 straff.} \\
\texttt{Vänligen försök igen!} \\
            }
        }
    }
    & \\ \\

    \theutterance \stepcounter{utterance}  
    & & & \multicolumn{4}{p{0.6\linewidth}}{
        \cellcolor[rgb]{0.9,0.9,0.9}{
            \makecell[{{p{\linewidth}}}]{
                \texttt{\tiny{[GM$\langle$P2]}}
                \texttt{SÄG: Både (3,6) och (6,6) har linjer hos mig (└ och ┘). Vad sägs om (3,7) eller (6,5)? Kan du kolla dem?} \\
            }
        }
    }
     \\ \\

    \theutterance \stepcounter{utterance}  
    & & & \multicolumn{2}{p{0.3\linewidth}}{
        \cellcolor[rgb]{0.9,0.9,0.9}{
            \makecell[{{p{\linewidth}}}]{
                \texttt{\tiny{[GM$|$GM]}}
                \texttt{SÄG: Både (3,6) och (6,6) har linjer hos mig (└ och ┘). Vad sägs om (3,7) eller (6,5)? Kan du kolla dem?} \\
            }
        }
    }
    & & \\ \\

    \theutterance \stepcounter{utterance}  
    & & \multicolumn{4}{p{0.6\linewidth}}{
        \cellcolor[rgb]{0.9,0.9,0.9}{
            \makecell[{{p{\linewidth}}}]{
                \texttt{\tiny{[P1$\langle$GM]}}
                \texttt{Ditt meddelande har vidarebefordrats till den andra spelaren.} \\
\\ 
\\ 
\texttt{Du spelar för närvarande runda 10 av maximalt 12.} \\
\texttt{Ni har tillsammans ackumulerat 8 av 8 straff.} \\
\texttt{Den andra spelaren skickade detta meddelande:} \\
\\ 
\\ 
\texttt{"Både (3,6) och (6,6) har linjer hos mig (└ och ┘). Vad sägs om (3,7) eller (6,5)? Kan du kolla dem?"} \\
\\ 
\\ 
\texttt{Vad är ditt nästa kommando?} \\
            }
        }
    }
    & \\ \\

    \theutterance \stepcounter{utterance}  
    & \multicolumn{4}{p{0.6\linewidth}}{
        \cellcolor[rgb]{0.9,0.9,0.9}{
            \makecell[{{p{\linewidth}}}]{
                \texttt{\tiny{[P1$\rangle$GM]}}
                \texttt{SÄG: Vi har tyvärr nått 8 straff och förlorat spelet. Men för att avsluta {-} (3,7) har en linje (│) hos mig, och (6,5) är tomt (◌). Det hade nog funkat! Bra försök i alla fall!} \\
            }
        }
    }
    & & \\ \\

    \theutterance \stepcounter{utterance}  
    & & & \multicolumn{2}{p{0.3\linewidth}}{
        \cellcolor[rgb]{0.9,0.9,0.9}{
            \makecell[{{p{\linewidth}}}]{
                \texttt{\tiny{[GM$|$GM]}}
                \texttt{SÄG: Vi har tyvärr nått 8 straff och förlorat spelet. Men för att avsluta {-} (3,7) har en linje (│) hos mig, och (6,5) är tomt (◌). Det hade nog funkat! Bra försök i alla fall!} \\
            }
        }
    }
    & & \\ \\

    \theutterance \stepcounter{utterance}  
    & & \multicolumn{4}{p{0.6\linewidth}}{
        \cellcolor[rgb]{0.9,0.9,0.9}{
            \makecell[{{p{\linewidth}}}]{
                \texttt{\tiny{[GM$\rangle$P2]}}
                \texttt{Ditt meddelande har vidarebefordrats till den andra spelaren.} \\
\\ 
\\ 
\texttt{Du spelar för närvarande runda 11 av maximalt 12.} \\
\texttt{Ni har tillsammans ackumulerat 8 av 8 straff.} \\
\texttt{Den andra spelaren skickade detta meddelande:} \\
\\ 
\\ 
\texttt{"Vi har tyvärr nått 8 straff och förlorat spelet. Men för att avsluta {-} (3,7) har en linje (│) hos mig, och (6,5) är tomt (◌). Det hade nog funkat! Bra försök i alla fall!"} \\
\\ 
\\ 
\texttt{Vad är ditt nästa kommando?} \\
            }
        }
    }
    & \\ \\

    \theutterance \stepcounter{utterance}  
    & & & \multicolumn{4}{p{0.6\linewidth}}{
        \cellcolor[rgb]{0.9,0.9,0.9}{
            \makecell[{{p{\linewidth}}}]{
                \texttt{\tiny{[GM$\langle$P2]}}
                \texttt{SÄG: färdig?} \\
            }
        }
    }
     \\ \\

    \theutterance \stepcounter{utterance}  
    & & & \multicolumn{2}{p{0.3\linewidth}}{
        \cellcolor[rgb]{0.9,0.9,0.9}{
            \makecell[{{p{\linewidth}}}]{
                \texttt{\tiny{[GM$|$GM]}}
                \texttt{SÄG: färdig?} \\
            }
        }
    }
    & & \\ \\

    \theutterance \stepcounter{utterance}  
    & & \multicolumn{4}{p{0.6\linewidth}}{
        \cellcolor[rgb]{0.9,0.9,0.9}{
            \makecell[{{p{\linewidth}}}]{
                \texttt{\tiny{[P1$\langle$GM]}}
                \texttt{Ditt meddelande har vidarebefordrats till den andra spelaren.} \\
\\ 
\\ 
\texttt{Du spelar för närvarande runda 11 av maximalt 12.} \\
\texttt{Ni har tillsammans ackumulerat 8 av 8 straff.} \\
\texttt{Den andra spelaren skickade detta meddelande:} \\
\\ 
\\ 
\texttt{"färdig?"} \\
\\ 
\\ 
\texttt{Vad är ditt nästa kommando?} \\
            }
        }
    }
    & \\ \\

    \theutterance \stepcounter{utterance}  
    & \multicolumn{4}{p{0.6\linewidth}}{
        \cellcolor[rgb]{0.9,0.9,0.9}{
            \makecell[{{p{\linewidth}}}]{
                \texttt{\tiny{[P1$\rangle$GM]}}
                \texttt{SÄG: färdig!} \\
            }
        }
    }
    & & \\ \\

    \theutterance \stepcounter{utterance}  
    & & & \multicolumn{2}{p{0.3\linewidth}}{
        \cellcolor[rgb]{0.9,0.9,0.9}{
            \makecell[{{p{\linewidth}}}]{
                \texttt{\tiny{[GM$|$GM]}}
                \texttt{SÄG: färdig!} \\
            }
        }
    }
    & & \\ \\

    \theutterance \stepcounter{utterance}  
    & & & \multicolumn{2}{p{0.3\linewidth}}{
        \cellcolor[rgb]{0.9,0.9,0.9}{
            \makecell[{{p{\linewidth}}}]{
                \texttt{\tiny{[GM$|$GM]}}
                \texttt{true} \\
            }
        }
    }
    & & \\ \\

    \theutterance \stepcounter{utterance}  
    & & & \multicolumn{2}{p{0.3\linewidth}}{
        \cellcolor[rgb]{0.9,0.9,0.9}{
            \makecell[{{p{\linewidth}}}]{
                \texttt{\tiny{[GM$|$GM]}}
                \texttt{Initial states:} \\
\texttt{```} \\
\texttt{Player 1:} \\
\texttt{ 1234567} \\
\texttt{╔══╤═╤══╗} \\
\texttt{║◌◌│◌│◌◌║ 1} \\
\texttt{║◌◌└─┼──╢ 2} \\
\texttt{║◌◌◌◌│P◌║ 3} \\
\texttt{║◌┌─┬┴┬─╢ 4} \\
\texttt{╟─┤◌│◌│◌║ 5} \\
\texttt{║◌└┬┴┬┘L║ 6} \\
\texttt{║◌◌│C│◌◌║ 7} \\
\texttt{╚══╧═╧══╝} \\
\\ 
\\ 
\texttt{Player 2:} \\
\texttt{ 1234567} \\
\texttt{╔══╤═╤══╗} \\
\texttt{║◌P│◌│◌◌║ 1} \\
\texttt{║◌L└─┼──╢ 2} \\
\texttt{║◌◌◌◌│◌◌║ 3} \\
\texttt{║◌┌─┬┴┬─╢ 4} \\
\texttt{╟─┤◌│◌│◌║ 5} \\
\texttt{║◌└┬┴┬┘◌║ 6} \\
\texttt{║◌◌│◌│◌C║ 7} \\
\texttt{╚══╧═╧══╝} \\
\\ 
\texttt{```} \\
            }
        }
    }
    & & \\ \\

    \theutterance \stepcounter{utterance}  
    & & & \multicolumn{2}{p{0.3\linewidth}}{
        \cellcolor[rgb]{0.9,0.9,0.9}{
            \makecell[{{p{\linewidth}}}]{
                \texttt{\tiny{[GM$|$GM]}}
                \texttt{End states:} \\
\texttt{```} \\
\texttt{Player 1:} \\
\texttt{ 1234567} \\
\texttt{╔══╤═╤══╗} \\
\texttt{║◌◌│◌│◌◌║ 1} \\
\texttt{║◌◌└─┼──╢ 2} \\
\texttt{║◌◌◌◌│P◌║ 3} \\
\texttt{║◌┌─┬┴┬─╢ 4} \\
\texttt{╟─┤◌│L│◌║ 5} \\
\texttt{║◌└┬┴┬┘◌║ 6} \\
\texttt{║◌◌│C│◌◌║ 7} \\
\texttt{╚══╧═╧══╝} \\
\\ 
\\ 
\texttt{Player 2:} \\
\texttt{ 1234567} \\
\texttt{╔══╤═╤══╗} \\
\texttt{║◌P│◌│◌◌║ 1} \\
\texttt{║◌◌└─┼──╢ 2} \\
\texttt{║◌◌◌◌│◌◌║ 3} \\
\texttt{║◌┌─┬┴┬─╢ 4} \\
\texttt{╟─┤◌│L│◌║ 5} \\
\texttt{║◌└┬┴┬┘◌║ 6} \\
\texttt{║◌◌│◌│◌C║ 7} \\
\texttt{╚══╧═╧══╝} \\
\\ 
\texttt{```} \\
            }
        }
    }
    & & \\ \\

    \theutterance \stepcounter{utterance}  
    & & & \multicolumn{2}{p{0.3\linewidth}}{
        \cellcolor[rgb]{0.9,0.9,0.9}{
            \makecell[{{p{\linewidth}}}]{
                \texttt{\tiny{[GM$|$GM]}}
                \texttt{* success: True} \\
\texttt{* lose: False} \\
\texttt{* aborted: False} \\
\texttt{{-}{-}{-}{-}{-}{-}{-}} \\
\texttt{* Shifts: 0.00} \\
\texttt{* Max Shifts: 4.00} \\
\texttt{* Min Shifts: 2.00} \\
\texttt{* End Distance Sum: 7.47} \\
\texttt{* Init Distance Sum: 13.88} \\
\texttt{* Expected Distance Sum: 12.57} \\
\texttt{* Penalties: 8.00} \\
\texttt{* Max Penalties: 8.00} \\
\texttt{* Rounds: 11.00} \\
\texttt{* Max Rounds: 12.00} \\
\texttt{* Object Count: 3.00} \\
            }
        }
    }
    & & \\ \\

\end{supertabular}
}

\end{document}
